\newpage
\phantomsection
\label{prf:interval-iff-order-convex}
\LRAProofFor{thm:interval-iff-order-convex}

\begin{remark*}[Return]
\hyperref[thm:interval-iff-order-convex]{Return to Theorem}
\end{remark*}

\begin{theorem*}[Intervals Are Exactly the Order-Convex Sets]
A nonempty $I\subseteq\mathbb{R}$ is an interval if and only if it is order-convex.
\end{theorem*}

\textbf{Professional Standard Proof.}~\\
\begin{proof}
TODO
\end{proof}

\textbf{Detailed Learning Proof.}~\\
TODO: Expand the professional proof into a detailed learning proof.

\begin{remark*}[Proof structure]
The forward direction is a direct check on each interval shape. The reverse direction uses completeness: set $a=\inf I$ (or $-\infty$) and $b=\sup I$ (or $+\infty$); order-convexity forces $(a,b)\subseteq I\subseteq[a,b]$, and the four endpoint cases give the precise shape.
\end{remark*}

\begin{dependencies}
\begin{itemize}
  \item \hyperref[def:order-convex-set]{Order-Convex Set}
  \item TODO: link to the Least Upper Bound Property.
\end{itemize}
\end{dependencies}
